%!TEX root = ../template.tex
%%%%%%%%%%%%%%%%%%%%%%%%%%%%%%%%%%%%%%%%%%%%%%%%%%%%%%%%%%%%%%%%%%%%
%% abstract-en.tex
%% NOVA thesis document file
%%
%% Abstract in English
%%%%%%%%%%%%%%%%%%%%%%%%%%%%%%%%%%%%%%%%%%%%%%%%%%%%%%%%%%%%%%%%%%%%

\typeout{NT FILE abstract-en.tex}%

Open-source software and commodity software-defined radios have made it possible to deploy a private 5G Standalone network without a commercial operator's budget. Deploying it, however, is only half the problem: a working network can still perform far below what its hardware allows, and tuning it is a challenge of its own.

This dissertation implemented a 5G Standalone network alongside a 4G reference network, sharing the same Open5GS core: the srsRAN Project gNB with a USRP-2901, and srsRAN 4G with a bladeRF xA4. The 5G network was evaluated at three distances over three rounds of tuning.

Against this reference, the untuned 5G uplink was more than 20 times slower at 10\,m, partly because a 3:1 TDD pattern limited its airtime. After rebalancing the pattern to 2:2, the uplink reached the 4G level, and lowering the receiver gain removed a short-range overload, reaching 28.5\,Mbps at 2\,m, twice the 4G reference, at the cost of range. The idle round-trip time averaged 25 to 28\,ms, but rose to seconds under TCP load due to an oversized default RLC queue.

These results show that an open-source 5G Standalone network can be deployed and operated as a laboratory network: once adapted to its hardware and environment, it matched the 4G reference and exceeded it at short range, and much of what limited it came from default or inherited parameters rather than from the software. Both networks remain available for future experimental work.
% Abstract keywords
\keywords{
  5G Standalone \and
  Open-source \and
  Software-Defined radio \and
  5G Testbed \and
  Performance evaluation
}
